% !TeX root = ../main.tex
\section{Структурная схема устройства}

Цифровой квадратурный приемный тракт выполняет оцифровку входного колебания, перенос спектра на промежуточную частоту с формированием квадратурных составляющих, децимацию и избирательную фильтрацию. Базовая структурная схема тракта приведена на рисунке \ref{fig:sch_main}.

\begin{figure}[H]
    \centering
    \includegraphics[width=0.85\textwidth]{schmain.png}
    \caption{Структурная схема квадратурного цифрового приемника}
    \label{fig:sch_main}
\end{figure}

Обработка сигналов включает следующие этапы:
\begin{enumerate}
    \item Оцифровка входной смеси в аналого-цифровом преобразователе (ADC) с разрядностью $noba = 13$ бит на частоте дискретизации $f_s = 180$ МГц.
    \item Квадратурное перемножение отсчетов с ортогональными гармоническими колебаниями цифрового гетеродина (NCO) несущей частоты $f_G = 210.0$ МГц и разрядностью $nobg = 13$ бит для разделения на синфазный ($I$) и квадратурный ($Q$) каналы.
    \item Первичная низкочастотная фильтрация и прореживание отсчетов в 10 раз ($R = 10$) с помощью каскадного интегрально-гребенчатого фильтра (CIC) 8-го порядка ($N = 8$).
    \item Компенсация частотных искажений в полосе пропускания и подавление внеполосных компонент с помощью корректирующего фильтра с конечной импульсной характеристикой (FIR).
\end{enumerate}


\subsection{Квадратурный перенос частоты}

На вход квадратурного смесителя поступает дискретизированный вещественный сигнал $x[n] = x_{\text{АЦП}}[n]$. Цифровой гетеродин формирует дискретные опорные колебания:
\begin{align}
    I[n] &= x[n] \cdot \cos(2\pi f_G n T_s), \label{eq:i_mix} \\
    Q[n] &= -x[n] \cdot \sin(2\pi f_G n T_s), \label{eq:q_mix}
\end{align}
где:
\begin{itemize}
    \item $T_s = 1/f_s$ --- период дискретизации входного тракта ($1/180$ МГц);
    \item $f_G = 210.0$ МГц --- частота цифрового гетеродина;
    \item $n$ --- целочисленный индекс отсчета времени.
\end{itemize}

В результате перемножения спектр полезного радиосигнала переносится на разностную промежуточную частоту $f_{IF} = f_c - f_G = 0.1$ МГц. Суммарная составляющая с частотой $f_c + f_G = 420.1$ МГц в результате дискретизации и периодичности спектра отображается в первую зону Найквиста на частоту $f_{\text{сум}} = -60.1$ МГц.

\subsection{Интегрально-гребенчатый дециматор (CIC)}

Каскадный фильтр Хогенауэра реализует децимацию без умножителей. Фильтр $N$-го порядка состоит из $N$ цифровых интеграторов, работающих на частоте $f_s$, узла прореживания на $R = 10$ и $N$ гребенчатых секций, работающих на частоте $f_{\text{дец}} = f_s / R = 18$ МГц.

\begin{figure}[H]
    \centering
    \includegraphics[width=0.7\textwidth]{CIC1.png}
    \caption{Структурная схема каскада цифровых интеграторов ($N = 8$)}
    \label{fig:cic1}
\end{figure}

Одиночный интегратор описывается разностным уравнением и передаточной функцией:
\begin{equation}
    y_I[n] = y_I[n-1] + x_I[n] \quad \Longleftrightarrow \quad H_I(z) = \frac{1}{1 - z^{-1}}.
\end{equation}

\begin{figure}[H]
    \centering
    \includegraphics[width=0.7\textwidth]{CIC2.png}
    \caption{Структурная схема каскада гребенчатых звеньев ($N = 8$)}
    \label{fig:cic2}
\end{figure}

Каждое гребенчатое звено выполняет дифференцирование на пониженной частоте отсчетов:
\begin{equation}
    y_C[m] = x_C[m] - x_C[m-D] \quad \Longleftrightarrow \quad H_C(z) = 1 - z^{-D},
\end{equation}
где $D = 1$ --- коэффициент дифференциальной задержки, $m$ --- индекс отсчетов на частоте децимации $f_{\text{дец}}$.

Результирующая передаточная функция CIC-дециматора 8-го порядка ($N = 8$) имеет вид:
\begin{equation}
    H_{\text{CIC}}(z) = \left( \frac{1 - z^{-R D}}{1 - z^{-1}} \right)^N = \left( \frac{1 - z^{-10}}{1 - z^{-1}} \right)^8.
    \label{eq:cic_tf}
\end{equation}

\subsection{Корректирующий КИХ-фильтр (FIR)}

Для устранения спада АЧХ фильтра Хогенауэра в полосе полезного сигнала и обеспечения затухания не менее 60 дБ в полосе задерживания применяется КИХ-фильтр прямой трансверсальной формы (рис. \ref{fig:fir}).

\begin{figure}[H]
    \centering
    \includegraphics[width=0.85\textwidth]{FIR.png}
    \caption{Трансверсальная структура корректирующего КИХ-фильтра}
    \label{fig:fir}
\end{figure}

Работа КИХ-фильтра описывается дискретной сверткой:
\begin{equation}
    y[m] = \sum_{k=0}^{M} h[k] \cdot x[m-k],
    \label{eq:fir_conv}
\end{equation}
где:
\begin{itemize}
    \item $h[k]$ --- коэффициенты импульсной характеристики фильтра;
    \item $M$ --- порядок КИХ-фильтра;
    \item $x[m]$ и $y[m]$ --- входная и выходная последовательности отсчетов на частоте $f_{\text{дец}} = 18$ МГц.
\end{itemize}
